\IfFileExists{stacks-project.cls}{%
\documentclass{stacks-project}
}{%
\documentclass{amsart}
}

% For dealing with references we use the comment environment
\usepackage{verbatim}
\newenvironment{reference}{\comment}{\endcomment}
%\newenvironment{reference}{}{}
\newenvironment{slogan}{\comment}{\endcomment}
\newenvironment{history}{\comment}{\endcomment}

% For commutative diagrams we use Xy-pic
\usepackage[all]{xy}

% We use 2cell for 2-commutative diagrams.
\xyoption{2cell}
\UseAllTwocells

% We use multicol for the list of chapters between chapters
\usepackage{multicol}

% This is generally recommended for better output
\usepackage{lmodern}
\usepackage[T1]{fontenc}

% For cross-file-references
\usepackage{xr-hyper}

% Package for hypertext links:
\usepackage{hyperref}

% For any local file, say "hello.tex" you want to link to please
% use \externaldocument[hello-]{hello}
\externaldocument[introduction-]{introduction}
\externaldocument[conventions-]{conventions}
\externaldocument[sets-]{sets}
\externaldocument[categories-]{categories}
\externaldocument[topology-]{topology}
\externaldocument[sheaves-]{sheaves}
\externaldocument[sites-]{sites}
\externaldocument[stacks-]{stacks}
\externaldocument[fields-]{fields}
\externaldocument[algebra-]{algebra}
\externaldocument[brauer-]{brauer}
\externaldocument[homology-]{homology}
\externaldocument[derived-]{derived}
\externaldocument[simplicial-]{simplicial}
\externaldocument[more-algebra-]{more-algebra}
\externaldocument[smoothing-]{smoothing}
\externaldocument[modules-]{modules}
\externaldocument[sites-modules-]{sites-modules}
\externaldocument[injectives-]{injectives}
\externaldocument[cohomology-]{cohomology}
\externaldocument[sites-cohomology-]{sites-cohomology}
\externaldocument[dga-]{dga}
\externaldocument[dpa-]{dpa}
\externaldocument[sdga-]{sdga}
\externaldocument[hypercovering-]{hypercovering}
\externaldocument[schemes-]{schemes}
\externaldocument[constructions-]{constructions}
\externaldocument[properties-]{properties}
\externaldocument[morphisms-]{morphisms}
\externaldocument[coherent-]{coherent}
\externaldocument[divisors-]{divisors}
\externaldocument[limits-]{limits}
\externaldocument[varieties-]{varieties}
\externaldocument[topologies-]{topologies}
\externaldocument[descent-]{descent}
\externaldocument[perfect-]{perfect}
\externaldocument[more-morphisms-]{more-morphisms}
\externaldocument[flat-]{flat}
\externaldocument[groupoids-]{groupoids}
\externaldocument[more-groupoids-]{more-groupoids}
\externaldocument[etale-]{etale}
\externaldocument[chow-]{chow}
\externaldocument[intersection-]{intersection}
\externaldocument[pic-]{pic}
\externaldocument[weil-]{weil}
\externaldocument[adequate-]{adequate}
\externaldocument[dualizing-]{dualizing}
\externaldocument[duality-]{duality}
\externaldocument[discriminant-]{discriminant}
\externaldocument[derham-]{derham}
\externaldocument[local-cohomology-]{local-cohomology}
\externaldocument[algebraization-]{algebraization}
\externaldocument[curves-]{curves}
\externaldocument[resolve-]{resolve}
\externaldocument[models-]{models}
\externaldocument[functors-]{functors}
\externaldocument[equiv-]{equiv}
\externaldocument[pione-]{pione}
\externaldocument[etale-cohomology-]{etale-cohomology}
\externaldocument[proetale-]{proetale}
\externaldocument[relative-cycles-]{relative-cycles}
\externaldocument[more-etale-]{more-etale}
\externaldocument[trace-]{trace}
\externaldocument[crystalline-]{crystalline}
\externaldocument[spaces-]{spaces}
\externaldocument[spaces-properties-]{spaces-properties}
\externaldocument[spaces-morphisms-]{spaces-morphisms}
\externaldocument[decent-spaces-]{decent-spaces}
\externaldocument[spaces-cohomology-]{spaces-cohomology}
\externaldocument[spaces-limits-]{spaces-limits}
\externaldocument[spaces-divisors-]{spaces-divisors}
\externaldocument[spaces-over-fields-]{spaces-over-fields}
\externaldocument[spaces-topologies-]{spaces-topologies}
\externaldocument[spaces-descent-]{spaces-descent}
\externaldocument[spaces-perfect-]{spaces-perfect}
\externaldocument[spaces-more-morphisms-]{spaces-more-morphisms}
\externaldocument[spaces-flat-]{spaces-flat}
\externaldocument[spaces-groupoids-]{spaces-groupoids}
\externaldocument[spaces-more-groupoids-]{spaces-more-groupoids}
\externaldocument[bootstrap-]{bootstrap}
\externaldocument[spaces-pushouts-]{spaces-pushouts}
\externaldocument[spaces-chow-]{spaces-chow}
\externaldocument[groupoids-quotients-]{groupoids-quotients}
\externaldocument[spaces-more-cohomology-]{spaces-more-cohomology}
\externaldocument[spaces-simplicial-]{spaces-simplicial}
\externaldocument[spaces-duality-]{spaces-duality}
\externaldocument[formal-spaces-]{formal-spaces}
\externaldocument[restricted-]{restricted}
\externaldocument[spaces-resolve-]{spaces-resolve}
\externaldocument[formal-defos-]{formal-defos}
\externaldocument[defos-]{defos}
\externaldocument[cotangent-]{cotangent}
\externaldocument[examples-defos-]{examples-defos}
\externaldocument[algebraic-]{algebraic}
\externaldocument[examples-stacks-]{examples-stacks}
\externaldocument[stacks-sheaves-]{stacks-sheaves}
\externaldocument[criteria-]{criteria}
\externaldocument[artin-]{artin}
\externaldocument[quot-]{quot}
\externaldocument[stacks-properties-]{stacks-properties}
\externaldocument[stacks-morphisms-]{stacks-morphisms}
\externaldocument[stacks-limits-]{stacks-limits}
\externaldocument[stacks-cohomology-]{stacks-cohomology}
\externaldocument[stacks-perfect-]{stacks-perfect}
\externaldocument[stacks-introduction-]{stacks-introduction}
\externaldocument[stacks-more-morphisms-]{stacks-more-morphisms}
\externaldocument[stacks-geometry-]{stacks-geometry}
\externaldocument[moduli-]{moduli}
\externaldocument[moduli-curves-]{moduli-curves}
\externaldocument[examples-]{examples}
\externaldocument[exercises-]{exercises}
\externaldocument[guide-]{guide}
\externaldocument[desirables-]{desirables}
\externaldocument[coding-]{coding}
\externaldocument[obsolete-]{obsolete}
\externaldocument[fdl-]{fdl}
\externaldocument[index-]{index}

% Theorem environments.
%
\theoremstyle{plain}
\newtheorem{theorem}[subsection]{Theorem}
\newtheorem{proposition}[subsection]{Proposition}
\newtheorem{lemma}[subsection]{Lemma}

\theoremstyle{definition}
\newtheorem{definition}[subsection]{Definition}
\newtheorem{example}[subsection]{Example}
\newtheorem{exercise}[subsection]{Exercise}
\newtheorem{situation}[subsection]{Situation}

\theoremstyle{remark}
\newtheorem{remark}[subsection]{Remark}
\newtheorem{remarks}[subsection]{Remarks}

\numberwithin{equation}{subsection}

% Macros
%
\def\lim{\mathop{\mathrm{lim}}\nolimits}
\def\colim{\mathop{\mathrm{colim}}\nolimits}
\def\Spec{\mathop{\mathrm{Spec}}}
\def\Hom{\mathop{\mathrm{Hom}}\nolimits}
\def\Ext{\mathop{\mathrm{Ext}}\nolimits}
\def\SheafHom{\mathop{\mathcal{H}\!\mathit{om}}\nolimits}
\def\SheafExt{\mathop{\mathcal{E}\!\mathit{xt}}\nolimits}
\def\Sch{\mathit{Sch}}
\def\Mor{\mathop{\mathrm{Mor}}\nolimits}
\def\Ob{\mathop{\mathrm{Ob}}\nolimits}
\def\Sh{\mathop{\mathit{Sh}}\nolimits}
\def\NL{\mathop{N\!L}\nolimits}
\def\CH{\mathop{\mathrm{CH}}\nolimits}
\def\proetale{{pro\text{-}\acute{e}tale}}
\def\etale{{\acute{e}tale}}
\def\QCoh{\mathit{QCoh}}
\def\Ker{\mathop{\mathrm{Ker}}}
\def\Im{\mathop{\mathrm{Im}}}
\def\Coker{\mathop{\mathrm{Coker}}}
\def\Coim{\mathop{\mathrm{Coim}}}

% Boxtimes
%
\DeclareMathSymbol{\boxtimes}{\mathbin}{AMSa}{"02}

%
% Macros for moduli stacks/spaces
%
\def\QCohstack{\mathcal{QC}\!\mathit{oh}}
\def\Cohstack{\mathcal{C}\!\mathit{oh}}
\def\Spacesstack{\mathcal{S}\!\mathit{paces}}
\def\Quotfunctor{\mathrm{Quot}}
\def\Hilbfunctor{\mathrm{Hilb}}
\def\Curvesstack{\mathcal{C}\!\mathit{urves}}
\def\Polarizedstack{\mathcal{P}\!\mathit{olarized}}
\def\Complexesstack{\mathcal{C}\!\mathit{omplexes}}
% \Pic is the operator that assigns to X its picard group, usage \Pic(X)
% \Picardstack_{X/B} denotes the Picard stack of X over B
% \Picardfunctor_{X/B} denotes the Picard functor of X over B
\def\Pic{\mathop{\mathrm{Pic}}\nolimits}
\def\Picardstack{\mathcal{P}\!\mathit{ic}}
\def\Picardfunctor{\mathrm{Pic}}
\def\Deformationcategory{\mathcal{D}\!\mathit{ef}}
\usepackage{cleveref}
\crefname{originalreference}{original reference}{original references}
\Crefname{originalreference}{Original reference}{Original references}
\makeatletter
\newlabel{algebra-lemma-exists-dvr}{{00PH}{source}{}{sourcepointer-1}{}}
\newlabel{algebra-lemma-exists-dvr@cref}{{[originalreference][0][]00PH}{[1][0][]source}{}{}{}}
\makeatother
\crefname{originaldefinition}{original definition}{original definitions}
\Crefname{originaldefinition}{Original definition}{Original definitions}
\makeatletter
\newlabel{definition-rig-smooth}{{0GCP}{source}{}{sourcepointer-2}{}}
\newlabel{definition-rig-smooth@cref}{{[originaldefinition][0][]0GCP}{[1][0][]source}{}{}{}}
\makeatother
\crefname{originalreference}{original reference}{original references}
\Crefname{originalreference}{Original reference}{Original references}
\makeatletter
\newlabel{formal-spaces-definition-formal-algebraic-space}{{0AIM}{source}{}{sourcepointer-3}{}}
\newlabel{formal-spaces-definition-formal-algebraic-space@cref}{{[originalreference][0][]0AIM}{[1][0][]source}{}{}{}}
\makeatother
\crefname{originallemma}{original lemma}{original lemmas}
\Crefname{originallemma}{Original lemma}{Original lemmas}
\makeatletter
\newlabel{lemma-base-change-rig-smooth}{{0GCR}{source}{}{sourcepointer-4}{}}
\newlabel{lemma-base-change-rig-smooth@cref}{{[originallemma][0][]0GCR}{[1][0][]source}{}{}{}}
\makeatother
\crefname{originallemma}{original lemma}{original lemmas}
\Crefname{originallemma}{Original lemma}{Original lemmas}
\makeatletter
\newlabel{lemma-base-change-rig-surjective}{{0AQS}{source}{}{sourcepointer-5}{}}
\newlabel{lemma-base-change-rig-surjective@cref}{{[originallemma][0][]0AQS}{[1][0][]source}{}{}{}}
\makeatother
\crefname{originallemma}{original lemma}{original lemmas}
\Crefname{originallemma}{Original lemma}{Original lemmas}
\makeatletter
\newlabel{lemma-equivalent-with-artin-smooth}{{0GAJ}{source}{}{sourcepointer-6}{}}
\newlabel{lemma-equivalent-with-artin-smooth@cref}{{[originallemma][0][]0GAJ}{[1][0][]source}{}{}{}}
\makeatother
\crefname{originallemma}{original lemma}{original lemmas}
\Crefname{originallemma}{Original lemma}{Original lemmas}
\makeatletter
\newlabel{lemma-etale-covering-rig-surjective}{{0AQU}{source}{}{sourcepointer-7}{}}
\newlabel{lemma-etale-covering-rig-surjective@cref}{{[originallemma][0][]0AQU}{[1][0][]source}{}{}{}}
\makeatother
\crefname{originallemma}{original lemma}{original lemmas}
\Crefname{originallemma}{Original lemma}{Original lemmas}
\makeatletter
\newlabel{lemma-exact-sequence-NL}{{0ALM}{source}{}{sourcepointer-8}{}}
\newlabel{lemma-exact-sequence-NL@cref}{{[originallemma][0][]0ALM}{[1][0][]source}{}{}{}}
\makeatother
\crefname{originallemma}{original lemma}{original lemmas}
\Crefname{originallemma}{Original lemma}{Original lemmas}
\makeatletter
\newlabel{lemma-exact-sequence-NL-rig-smooth}{{0GH2}{source}{}{sourcepointer-9}{}}
\newlabel{lemma-exact-sequence-NL-rig-smooth@cref}{{[originallemma][0][]0GH2}{[1][0][]source}{}{}{}}
\makeatother
\crefname{originallemma}{original lemma}{original lemmas}
\Crefname{originallemma}{Original lemma}{Original lemmas}
\makeatletter
\newlabel{lemma-finite-type-morphisms}{{0GC0}{source}{}{sourcepointer-10}{}}
\newlabel{lemma-finite-type-morphisms@cref}{{[originallemma][0][]0GC0}{[1][0][]source}{}{}{}}
\makeatother
\crefname{originallemma}{original lemma}{original lemmas}
\Crefname{originallemma}{Original lemma}{Original lemmas}
\makeatletter
\newlabel{lemma-rig-smooth-rig-flat}{{0GH1}{source}{}{sourcepointer-11}{}}
\newlabel{lemma-rig-smooth-rig-flat@cref}{{[originallemma][0][]0GH1}{[1][0][]source}{}{}{}}
\makeatother
\title{Rig-smooth and rig-etale descent for morphisms of formal algebraic spaces}
\author{The Stacks Project Authors}
\date{stacks-project:restricted.tex:local-snapshot-20261012; GNU FDL 1.2 or later}
\begin{document}
\maketitle
\noindent\textit{Editorial note.} Counted span is the author's own proof environment, source lines 5638-5732, following the statement at 5625-5636.\par\medskip
\noindent\textit{Editorial note.} The 'some details omitted' at the contradiction setup is the instantiation of the cited characterization lemma, flagged in the notes; the rig-etale case is the mirror argument with both cohomology objects vanishing.\par\medskip
\noindent\textit{Editorial note.} Steps the author leaves to the reader are flagged here for the reviewer: source line(s) 5676 carry an "omitted"/"follows easily"/"immediate from the definitions" phrase. Those sentences are retained verbatim; the displayed argument around them stays inside the counted span.\par\medskip
\noindent\textit{Allowed standard prerequisites (editorial).}\begin{itemize}
\item Characterization of rig-smooth maps by the completed naive cotangent complex. lemma-equivalent-with-artin-smooth (same chapter, cited at \textasciitilde{}5671; its own proof has stated omissions and was not counted). Failure of rig-smoothness gives a prime with H\textasciicircum{}-1 nonzero or H\textasciicircum{}0 not finite free.
\item Flatness of rig-smooth ring maps. lemma-rig-smooth-rig-flat (carded L26159, cited at \textasciitilde{}5704). B\_q -> A\_p is flat, hence faithfully flat.
\item Exact sequence of completed naive cotangent complexes. lemma-exact-sequence-NL and lemma-exact-sequence-NL-rig-smooth (same chapter, cited at \textasciitilde{}5710). Transmits the vanishing and finite-free conclusions along C -> B -> A.
\item DVR lifting through etale coverings. Algebra, Lemma algebra-lemma-exists-dvr (cross-chapter) and lemma-etale-covering-rig-surjective (same chapter, cited at \textasciitilde{}5695). Produces the test ring R and lifts the Spf(R)-point to some Spf(A\_i).
\end{itemize}\medskip
\ifdefined\hypertarget\hypertarget{sourcepointer-1}{}\fi
\noindent\textit{Source reference.} Original reference 00PH: focused-seeds/geometry/S9-restricted-rig-etale-descent/ref-tags.txt, source lines 5--5. Navigation pointer only: the counted proof cites this statement by \textbackslash{}ref as a standing input; the deduction the counted proof has to make is carried out inside the counted span, and no part of that deduction is delegated to this pointer.\par\medskip
\ifdefined\hypertarget\hypertarget{sourcepointer-2}{}\fi
\noindent\textit{Source reference.} Original definition 0GCP: sources/stacks-restricted/pinned-original-source.tex, source lines 4622--4622. Navigation pointer only: the counted proof cites this statement by \textbackslash{}ref as a standing input; the deduction the counted proof has to make is carried out inside the counted span, and no part of that deduction is delegated to this pointer.\par\medskip
\ifdefined\hypertarget\hypertarget{sourcepointer-3}{}\fi
\noindent\textit{Source reference.} Original reference 0AIM: focused-seeds/geometry/S9-restricted-rig-etale-descent/ref-tags.txt, source lines 6--6. Navigation pointer only: the counted proof cites this statement by \textbackslash{}ref as a standing input; the deduction the counted proof has to make is carried out inside the counted span, and no part of that deduction is delegated to this pointer.\par\medskip
\ifdefined\hypertarget\hypertarget{sourcepointer-4}{}\fi
\noindent\textit{Source reference.} Original lemma 0GCR: sources/stacks-restricted/pinned-original-source.tex, source lines 4692--4692. Navigation pointer only: the counted proof cites this statement by \textbackslash{}ref as a standing input; the deduction the counted proof has to make is carried out inside the counted span, and no part of that deduction is delegated to this pointer.\par\medskip
\ifdefined\hypertarget\hypertarget{sourcepointer-5}{}\fi
\noindent\textit{Source reference.} Original lemma 0AQS: sources/stacks-restricted/pinned-original-source.tex, source lines 5324--5324. Navigation pointer only: the counted proof cites this statement by \textbackslash{}ref as a standing input; the deduction the counted proof has to make is carried out inside the counted span, and no part of that deduction is delegated to this pointer.\par\medskip
\ifdefined\hypertarget\hypertarget{sourcepointer-6}{}\fi
\noindent\textit{Source reference.} Original lemma 0GAJ: sources/stacks-restricted/pinned-original-source.tex, source lines 843--843. Navigation pointer only: the counted proof cites this statement by \textbackslash{}ref as a standing input; the deduction the counted proof has to make is carried out inside the counted span, and no part of that deduction is delegated to this pointer.\par\medskip
\ifdefined\hypertarget\hypertarget{sourcepointer-7}{}\fi
\noindent\textit{Source reference.} Original lemma 0AQU: sources/stacks-restricted/pinned-original-source.tex, source lines 5363--5363. Navigation pointer only: the counted proof cites this statement by \textbackslash{}ref as a standing input; the deduction the counted proof has to make is carried out inside the counted span, and no part of that deduction is delegated to this pointer.\par\medskip
\ifdefined\hypertarget\hypertarget{sourcepointer-8}{}\fi
\noindent\textit{Source reference.} Original lemma 0ALM: sources/stacks-restricted/pinned-original-source.tex, source lines 725--725. Navigation pointer only: the counted proof cites this statement by \textbackslash{}ref as a standing input; the deduction the counted proof has to make is carried out inside the counted span, and no part of that deduction is delegated to this pointer.\par\medskip
\ifdefined\hypertarget\hypertarget{sourcepointer-9}{}\fi
\noindent\textit{Source reference.} Original lemma 0GH2: sources/stacks-restricted/pinned-original-source.tex, source lines 4561--4561. Navigation pointer only: the counted proof cites this statement by \textbackslash{}ref as a standing input; the deduction the counted proof has to make is carried out inside the counted span, and no part of that deduction is delegated to this pointer.\par\medskip
\ifdefined\hypertarget\hypertarget{sourcepointer-10}{}\fi
\noindent\textit{Source reference.} Original lemma 0GC0: sources/stacks-restricted/pinned-original-source.tex, source lines 2556--2556. Navigation pointer only: the counted proof cites this statement by \textbackslash{}ref as a standing input; the deduction the counted proof has to make is carried out inside the counted span, and no part of that deduction is delegated to this pointer.\par\medskip
\ifdefined\hypertarget\hypertarget{sourcepointer-11}{}\fi
\noindent\textit{Source reference.} Original lemma 0GH1: sources/stacks-restricted/pinned-original-source.tex, source lines 4460--4460. Navigation pointer only: the counted proof cites this statement by \textbackslash{}ref as a standing input; the deduction the counted proof has to make is carried out inside the counted span, and no part of that deduction is delegated to this pointer.\par\medskip
\begin{lemma}
\label{lemma-rig-etale-descent}
Let $S$ be a scheme. Let $f : X \to Y$ and $g : Y \to Z$
be morphisms of locally Noetherian formal algebraic spaces over $S$.
Assume
\begin{enumerate}
\item $g$ is locally of finite type,
\item $f$ is rig-smooth (resp.\ rig-\'etale) and rig-surjective,
\item $g \circ f$ is rig-smooth (resp.\ rig-\'etale)
\end{enumerate}
then $g$ is rig-smooth (resp.\ rig-\'etale).
\end{lemma}

\begin{proof}
We will prove this in the rig-smooth case and indicate the necessary
changes to prove the rig-\'etale case at the end of the proof.
Consider a commutative diagram
$$
\xymatrix{
X \times_Y V \ar[r] \ar[d] &
V \ar[d] \ar[r] &
W \ar[d] \\
X \ar[r] &
Y \ar[r] &
Z
}
$$
with $V$ and $W$ affine formal algebraic spaces, $V \to Y$ and $W \to Z$
representable by algebraic spaces and \'etale. We have to show that
$V \to W$ corresponds to a rig-smooth map of adic Noetherian topological
rings, see Definition \ref{definition-rig-smooth}.
We may write $V = \text{Spf}(B)$ and $W = \text{Spf}(C)$ and that
$V \to W$ corresponds to an adic ring map $C \to B$ which is topologically
of finite type, see Lemma \ref{lemma-finite-type-morphisms}.

\medskip\noindent
We will use below without further mention that $X \times_Y V \to V$
is rig-smooth and rig-surjective, see
Lemmas \ref{lemma-base-change-rig-smooth} and
\ref{lemma-base-change-rig-surjective}.
Also, the composition $X \times_Y V \to V \to W$ is rig-smooth
since $g \circ f$ is rig-smooth.

\medskip\noindent
Let $I \subset C$ be an ideal of definition. The module
Assume $C \to B$ is not rig-smooth to get a contradiction.
This means that there exists a prime ideal $\mathfrak q \subset B$
not containing $IB$ such that either $H^{-1}(\NL_{B/C}^\wedge)_\mathfrak p$
is nonzero or $H^0(\NL_{B/C}^\wedge)_\mathfrak p$ is not a finite free
$B_\mathfrak q$-module. See
Lemma \ref{lemma-equivalent-with-artin-smooth}; some details
omitted. We may choose a maximal ideal
$IB + \mathfrak q \subset \mathfrak m$. By
Algebra, Lemma \ref{algebra-lemma-exists-dvr}
we can find a complete discrete valuation ring $R$ and an injective
local ring homomorphism $(B/\mathfrak q)_\mathfrak m \to R$.

\medskip\noindent
After replacing $R$ by an extension, we may assume given a lift
$\text{Spf}(R) \to X \times_Y V$ of the adic morphism
$\text{Spf}(R) \to V = \text{Spf}(B)$. Choose an \'etale
covering $\{\text{Spf}(A_i) \to X \times_Y V\}$ as in
Formal Spaces, Definition
\ref{formal-spaces-definition-formal-algebraic-space}.
By Lemma \ref{lemma-etale-covering-rig-surjective}
we may assume $\text{Spf}(R) \to X \times_Y V$ lifts to a
morphism $\text{Spf}(R) \to \text{Spf}(A_i)$ for some $i$
(this might require replacing $R$ by another extension).
Set $A = A_i$. Consider the ring maps
$$
C \to B \to A \to R
$$
Let $\mathfrak p \subset A$ be the kernel of the map
$A \to R$ and note that $\mathfrak p$ lies over $\mathfrak q$.
We know that $C \to A$ and $B \to A$ are rig-smooth.
In particular the ring map $B_\mathfrak q \to A_\mathfrak p$
is flat by Lemma \ref{lemma-rig-smooth-rig-flat}.
Consider the associated exact sequence
$$
\xymatrix{
&
H^0(\NL_{B/C}^\wedge) \otimes_B A_\mathfrak p \ar[r] &
H^0(\NL_{A/C}^\wedge)_\mathfrak p \ar[r] &
H^0(\NL_{A/B}^\wedge)_\mathfrak p \ar[r] & 0 \\
0 \ar[r] &
H^{-1}(\NL_{B/C}^\wedge \otimes_B A)_\mathfrak p \ar[r] &
H^{-1}(\NL_{A/C}^\wedge)_\mathfrak p \ar[r] &
H^{-1}(\NL_{A/B}^\wedge)_\mathfrak p \ar[llu]
}
$$
of Lemmas \ref{lemma-exact-sequence-NL} and
\ref{lemma-exact-sequence-NL-rig-smooth}.
Given the rig-smoothness of $C \to A$ and $B \to A$
we conclude that $H^{-1}(\NL_{B/C}^\wedge \otimes_B A)_\mathfrak p = 0$
and that $H^0(\NL_{B/C}^\wedge) \otimes_B A_\mathfrak p$
is finite free as a kernel of a surjection of finite free
$A_\mathfrak p$-modules. Since $B_\mathfrak q \to A_\mathfrak p$
is flat and hence faithfully flat, this implies that
$H^{-1}(\NL_{B/C}^\wedge)_\mathfrak q = 0$
and that $H^0(\NL_{B/C}^\wedge)_\mathfrak q$
is finite free which is the contradiction we were looking for.

\medskip\noindent
In the rig-\'etale case one argues in exactly the same manner
but the conclusion obtained is that both
$H^{-1}(\NL_{B/C}^\wedge)_\mathfrak q$
and $H^0(\NL_{B/C}^\wedge)_\mathfrak q$ are zero.
\end{proof}

\end{document}
